% ===== Intervals (6 problems) =====
\section{Intervals \texorpdfstring{{\normalsize\color{muted}$\cdot$ 6 problems}}{- 6 problems}}
\markboth{Intervals}{}

\begin{problem}{57}{Insert Interval}{Medium}
\summary{Insert a new interval into a sorted, non-overlapping list and merge as needed.}
\keyidea{Three phases: copy intervals ending before it; merge all overlapping into it; copy the rest.}
\complexity{O(n) time, O(n) space}
\begin{code}
class Solution:
    def insert(self, intervals: List[List[int]], newInterval: List[int]) -> List[List[int]]:
        res, i, n = [], 0, len(intervals)
        while i < n and intervals[i][1] < newInterval[0]:
            res.append(intervals[i]); i += 1
        while i < n and intervals[i][0] <= newInterval[1]:
            newInterval = [min(newInterval[0], intervals[i][0]),
                           max(newInterval[1], intervals[i][1])]
            i += 1
        res.append(newInterval)
        res.extend(intervals[i:])
        return res
\end{code}
\end{problem}

\begin{problem}{56}{Merge Intervals}{Medium}
\summary{Merge all overlapping intervals.}
\keyidea{Sort by start; extend the last merged interval if the next one starts before it ends.}
\complexity{O(n log n) time, O(n) space}
\begin{code}
class Solution:
    def merge(self, intervals: List[List[int]]) -> List[List[int]]:
        intervals.sort(key=lambda x: x[0])
        res = [intervals[0][:]]
        for s, e in intervals[1:]:
            if s <= res[-1][1]:
                res[-1][1] = max(res[-1][1], e)
            else:
                res.append([s, e])
        return res
\end{code}
\end{problem}

\begin{problem}{435}{Non-overlapping Intervals}{Medium}
\summary{Minimum intervals to remove so the rest don't overlap.}
\keyidea{Sort by start; on overlap, remove the one with the larger end (keep the earlier-ending one).}
\complexity{O(n log n) time, O(1) space}
\begin{code}
class Solution:
    def eraseOverlapIntervals(self, intervals: List[List[int]]) -> int:
        intervals.sort()
        removed, prev_end = 0, intervals[0][1]
        for s, e in intervals[1:]:
            if s >= prev_end:
                prev_end = e
            else:
                removed += 1
                prev_end = min(prev_end, e)
        return removed
\end{code}
\end{problem}

\begin{problem}{252}{Meeting Rooms}{Easy}
\summary{Can one person attend all meetings (no overlaps)?}
\keyidea{Sort by start; any meeting starting before the previous ends is a conflict.}
\complexity{O(n log n) time, O(1) space}
\begin{code}
class Solution:
    def canAttendMeetings(self, intervals: List[List[int]]) -> bool:
        intervals.sort()
        for i in range(1, len(intervals)):
            if intervals[i][0] < intervals[i - 1][1]:
                return False
        return True
\end{code}
\end{problem}

\begin{problem}{253}{Meeting Rooms II}{Medium}
\summary{Minimum number of rooms needed for all meetings.}
\keyidea{Sort starts and ends separately; sweep: a start before the earliest end needs a new room.}
\complexity{O(n log n) time, O(n) space}
\begin{code}
class Solution:
    def minMeetingRooms(self, intervals: List[List[int]]) -> int:
        starts = sorted(i[0] for i in intervals)
        ends = sorted(i[1] for i in intervals)
        rooms = best = e = 0
        for s in starts:
            while e < len(ends) and ends[e] <= s:
                rooms -= 1
                e += 1
            rooms += 1
            best = max(best, rooms)
        return best
\end{code}
\end{problem}

\begin{problem}{1851}{Minimum Interval to Include Each Query}{Hard}
\summary{For each query point, return the size of the smallest interval containing it (-1 if none).}
\keyidea{Process queries sorted; push intervals whose start \textless{}= q into a min-heap by size; pop those ending before q.}
\complexity{O((n + q) log n) time, O(n + q) space}
\begin{code}
class Solution:
    def minInterval(self, intervals: List[List[int]], queries: List[int]) -> List[int]:
        intervals.sort()
        h, ans, i = [], {}, 0
        for q in sorted(queries):
            while i < len(intervals) and intervals[i][0] <= q:
                l, r = intervals[i]
                heapq.heappush(h, (r - l + 1, r))
                i += 1
            while h and h[0][1] < q:
                heapq.heappop(h)
            ans[q] = h[0][0] if h else -1
        return [ans[q] for q in queries]
\end{code}
\end{problem}
